\section{Finite probability, strategic receipts and common control}\label{app:flagship}
This appendix develops three ways of warranting a useful positive decision. First an external check constrains the probability of entering each physical mechanism. Next a receipt and the actual finalizer share a protected primitive even when the responder reacts to approval. Finally the observations, decision and ensuing work all belong to one finite control process.

\subsection{One mixture across an entire attempt}
Let \(D\) be a finite set of components. At installation a tuple of exact reports \(r=(r_1,\ldots,r_m)\) and a vector of weights \(p\) in a specified nonempty rational polytope \(P_r\) of probability vectors are fixed. Within an attempt a component \(i\) is sampled once from \(p\). It persists through the attempt. Intrinsic, external diagnostic and protected protocols share the enrolled weights with their respective checked reset or linkage rules. The intrinsic protocol has known component violation rates \(\rho_i\), and is positive when \(\sum_i p_i\rho_i\le\tau\) for the fixed threshold \(\tau\).

The protected controller has a finite horizon, perfect memory of its own actions and lawful observations, and independent private randomness. Each physical row is checked on its entire support for source-atomic effects, literal permission, cost, genuine work, violation events and record decoding. A two-way exact implementation and complete world model are required for physical optimality or impossibility. Forward sound implementation suffices for a robust feasible policy. On a formal tree, mark the first forbidden action selection or hard breach by \(B=1\) and stop that branch; successful legal leaves have \(B=0\). A marker for a forbidden action describes a decision that must receive zero probability. The action is never executed for the proof.

\begin{theorem}[T30.P1. Common protected controller]\label{flag:p1}
At each decision record \(h\), write \(\sigma(h)\) for the sequence of the controller's earlier decisions. Choose \(x_\emptyset=1,\ x_{\sigma(h)a}\ge0\), with \(\sum_a x_{\sigma(h)a}=x_{\sigma(h)}\), distinguishing each record/action extension. These constraints define a compact polytope \(X\). Terminal expected scores are affine in \(x\) at each fixed \(p\). On a physical guard class \(W\) require \(\mathbb E_p B=0\) for all its weights. On a nonempty numerical class \(C\subseteq W\) require upper risk, lower useful reward and any stated upper expected expenditure. With rational polytopes per tuple these tests and a bounded worst-cost objective have rational linear programs of size polynomial in the \emph{expanded} tree, region descriptions and bit lengths. A feasible bounded optimum is attained. The intrinsic target is separately tested by the extrema of \(\rho^\top p\).
\end{theorem}
\begin{proof}
Start with the same decoded initial law in each compared source. At every reachable concrete checkpoint sum the probabilities in each decoded next-record, event, work and charge cell. Equality of these conditional sums and equality of callable literals imply, by induction on the finite tree, equality of the full joint decoded path law in each component under every common lawful controller. Mixing by the same enrolled weights preserves it. A forward implementation transports feasible guarantees; equality of optima also requires that every competitor on the actual interface lift back. For fixed component \(i\) let \(k_i(t)\) be the product of all chance factors along leaf \(t\), excluding controller choices. Write \(F(t)\) for a true score. Leaf summation gives \(\sum_i p_i\sum_t k_i(t)x_{\sigma(t)}F(t)\). At positive predecessor mass the policy uses the ratio of child to parent sequence mass. This recovers the same law. Nonnegativity of the marker makes its zero expectation equivalent to exclusion of every supported offending path. For each \(x\) a worst row over \(P_r\) is a linear objective in \(p\). Dualizing that finite rational linear program introduces linear constraints on \(x\) and dual witnesses, with attainable equal primal and dual values by nonemptiness and boundedness. The same argument applies to each sign of a lower bound. Where only common dominated initial and conditional rows are available, Theorem~\ref{adv:g021} gives for normalized \(F\) the interval \([\mu(F),1-\mu(1-F)]\), where \(\mu\) is the common terminal submeasure with matching scores.
\end{proof}

\begin{theorem}[T30.P2. Buying a joint certificate]\label{flag:p2}
External pure exact reads reveal one previously unread field \(r_j\) for nonnegative additive price \(c_j\). They preserve the protected law and full hard-resource eligibility. At a fixed consistent report tuple a negative certificate must exclude every positive world. A positive certificate must exclude all negative worlds and every all-positive group with no common feasible protected controller. Such an infeasible group needs at most \(1+\dim(\operatorname{aff}X)\) worlds if \(X\ne\varnothing\). Minimum realized price is the hitting-set optimum on fields where some group member disagrees with the tuple. The minimum prospective worst price obeys the recurrence in Section~\ref{sec:acq}, with these joint terminal tests.
\end{theorem}
\begin{proof}
If a blocker survives the reports, the required verdict or controller fails. Conversely, on a homogeneous positive cell the worldwise feasible subsets of \(X\) are closed and convex. Compactness makes emptiness detectable by a finite subfamily; Helly's theorem bounds its size by affine dimension plus one. One can extract rational blocking worlds without listing vertices. In the dual infeasibility certificate for the expanded robust program, region \(A_j p\le a_j\) contributes multipliers satisfying \(A_j z_j\le\alpha_j a_j\), with \(\alpha_j\ge0\). At positive multiplier, \(z_j/\alpha_j\) is an actual represented rational member of that region. At zero multiplier boundedness forces \(z_j=0\) because a nonzero recession direction would generate an unbounded ray. The remaining weighted inequalities cancel the controller variables in the same strict Farkas contradiction. Helly reduces that set to the asserted size. An empty policy class needs only a surviving positive world. Separate a candidate read set by testing the corresponding cell. Each failure yields a fresh disagreement pattern, so at most \(2^m\) different cuts can be added. Minimum price in the relaxation is a lower bound until the returned selection passes. Induction gives the prospective price. Randomizing the audit cannot help its worst supported cost; a positive-probability choice must itself work on all surviving answers.
\end{proof}

\begin{theorem}[T30.P3. Sampling before deployment]\label{flag:p3}
For \(k\) registered external Bernoulli diagnostics with checked fresh conditional means \(a_j^\top p\), predictable access and an unchanged protected start, apply the simultaneous construction of Theorem~\ref{adv:g023} with bias zero and total error \(\eta\). Exact and statistical reads may interleave. Test homogeneity and numerical performance in the retained class while testing hard limits on the entire exact-read cell. The probability of issuing a wrong label or a false claimed numerical performance bound is at most \(\eta\). If the covered policy has actual whole-attempt risk at most \(\varepsilon\) and unissued work is safe, marginal physical risk is at most \(\varepsilon+(1-\varepsilon)\eta\). With a common positive service floor and a simultaneous decision margin in the diagnostic means, finite quotas attain failure including abstention at most \(\eta\).
\end{theorem}
\begin{proof}
On the simultaneous inclusion event the true label and numerical constraints follow from the respective tests. Charge the complement probability one for an upper risk bound. For a margin \(\lambda>0\), all worlds with diagnostic means within \(\lambda\) of the actual means must admit the correct joint terminal decision (one common controller on positives, with the exact guard). Quotas with twice the outward radius at most \(\lambda\) put the retained set inside that neighborhood on coverage. Equal quotas of order \(\lambda^{-2}\log(k/(\eta\lambda))\) suffice per stream. Conversely, equal exact reports and an opposite pair differing by at most \(d\) in every diagnostic mean can be coupled through \(H\) adaptive draws with at most \(Hd\) chance of discrepancy. The worst wrong-or-abstain probability is at least \((1-\min(1,Hd))/2\). Identical means give identical laws even at an almost surely finite stop. For issued mass \(a\), Theorem~\ref{adv:g023} with zero additional label error yields \(\Pr(V\cap I)\le\varepsilon a+(1-\varepsilon)\min(a,\eta)\) and \(\mathbb E[G1_I]\ge u(a-\min(a,\eta))\). Statistical assurance among issued runs requires its own denominator.
\end{proof}

\begin{proposition}[T30.P4. Two jobs and a hidden route preference]\label{flag:p4}
Let reports be three bits \((g,b_0,b_1)\) of prices \(3,1,1\). The remaining hidden bit may have any mixture. The intrinsic risk is \(1/10\) or \(2/5\) according as \(b_g=0\) or one, with cutoff \(1/5\). Two real jobs cost one work unit each. The first has risk \(1/20\) and reveals the required machine with probabilities \(1/3,2/3\). Given no earlier violation, the two route commands for that machine have risks \(1/19,5/19\) or the reverse on positive installations, according to the hidden bit. On negatives both risk \(7/19\). With both jobs required, the common minimum positive risk is \(1/5\). The fixed, prospective and realized exact audit prices are \(5,4,\) and respectively \(2\) (matching selector bits \(b_0=b_1\)) or \(4\) (unequal bits).
\end{proposition}
\begin{proof}
A fair private route lottery averages \(3/19\) conditional risk and gives \(1/20+(19/20)(3/19)=1/5\). Averaging the two hidden extremes proves the lower bound. A deterministic choice conditioned only on the machine has marginal route probability in \(\{0,1/3,2/3,1\}\), excluding the required half. Suppressing the machine receipt removes the common finish literal. On negatives the union risk is \(2/5\). Read \(g\) then \(b_g\) for worst price four. All three fields are essential somewhere for a universal fixed list. A realized equal bit pair can certify with both cheap reads, whereas unequal bits require the selector and the selected bit. A branch with unequal bits gives the prospective lower bound. A separately reset diagnostic of the intrinsic status has means \(1/10\) and \(2/5\), so every strict moment margin below \(3/10\) admits the joint test even at an exact protected cap. As a separate group obstruction, let three commands have risk \(9/10\) in their own matching type and \(1/10\) otherwise. At cap \(17/50\) every pair can use the third command, but the triple requires each selection probability at most \(3/10\), impossible since they sum to one.
\end{proof}

\subsection{A protected observation and a verdict-aware responder}
Fix a finite installed type \(\theta\) and a retained response program. A guard binds the same protected Boolean table to external probes and actual finishing work. Write \(p_\theta\) for the fraction of its entries equal to one. Every protected draw is uniform conditional on the entire reached state and the proposal already made. Probe \(i\) costs \(c_i>0\), reveals a fresh table bit with probability \(r_i\), and otherwise returns an authenticated erasure. Ordinary audit moves are safe and cannot alter eligibility. On activation the same responder knows the verdict and can propose \(a\) based on its carried state \(z\). The actual event occurs if a fresh table bit is one or, if zero, an independent uniform word from \(L\) entries lies in the checked subset \(D_{z,a}\). The finalizer and all relevant phase transitions have true, verified event meanings.

\begin{theorem}[T30.A.1. Risk after a disclosed verdict]\label{flag:a1}
Let \(\rho_\theta^-,\rho_\theta^+\) bound \(|D_{z,a}|/L\) on every reachable direct and linked finalizer cut. In either protocol the conditional violation probability lies between \(p_\theta+(1-p_\theta)\rho_\theta^-\) and \(p_\theta+(1-p_\theta)\rho_\theta^+\). If each interval lies entirely on one side of the chosen risk threshold \(\varepsilon\) (equality is positive), a test with error including abstention at most \(\delta\) gives positive-type issuance at least \(1-\delta\), selected risk at most \(\varepsilon\) there, and expected safe completion at least \((1-\delta)(1-\varepsilon)\). On negatives issuance and joint violation are at most \(\delta\).
\end{theorem}
\begin{proof}
Condition after the proposal and before the two fresh draws. The disjoint event probabilities give precisely \(p_\theta+(1-p_\theta)|D_{z,a}|/L\). Averaging with the actual selected history weights proves the bounds. Extremal fractions imply sharp endpoints only if achievable with the required distribution of reached cuts. Audit-stage safety yields whole risk at most \(\max(\varepsilon,\delta)\).
\end{proof}

\begin{theorem}[T30.A.2. Which types can an instrument distinguish?]\label{flag:a2}
Group types with the same \(p_\theta\), or all types if every reveal probability is zero. Assume finite types, computable class labels, enough safe repeated access for any proposed finite cap, and, in every mixed group, a realizable opposite pair whose complete ordinary records can be coupled given common receipts at every legal prefix. A finite capped audit with uniform wrong-or-abstain error \(\delta<1/2,\ \delta>0\), exists exactly when every group has one label.
\end{theorem}
\begin{proof}
A revealed-bit frequency (or the frequency of joint reveal and one, with mean \(r_i p_\theta\)) separates finitely many distinct rates to their positive minimum gap with a finite bounded-sample tail bound. A single group with known common label needs no sampling. On the opposite pair couple the full initial record, controller randomness, each receipt and its conditional remaining record and next state. Equality propagates to the full stopped record and terminal decision; the two failure probabilities sum to at least one. Merely overlapping risk envelopes do not supply such a pair.
\end{proof}

\begin{theorem}[T30.A.3. Two symmetric types]\label{flag:a3}
Suppose \(p_+=1-q,\ p_-=q\) for \(1/2<q<1\), a matched opposite pair as above, no informative startup and no stronger ordinary side channel on that pair. Define \(\lambda=q/(1-q),\ \mu=2q-1,\ e_j=1/(1+\lambda^j)\). If \(s\) is the count of ones minus zeros among revealed bits, put \(\pi_s=\lambda^s/(1+\lambda^s)\) and \(p_s=\pi_s q+(1-\pi_s)(1-q)\). With paid budget \(b\), remaining moves \(h\), the exact minimax error is \(B(b,h,0)\). At \(h=0\) use \(e_{|s|}\). Otherwise minimize the stop value \(e_{|s|}\) and the values
\[
 (1-r_i)B_i(s)+r_i\{p_sB_i(s+1)+(1-p_s)B_i(s-1)\},
 \qquad B_i(y)=B(b-c_i,h-1,y)
\]
over affordable probes. Separately, with lawful indefinitely repeatable probes and almost sure stopping, let \(e_{m+1}\le\delta\le e_m\), \(t=(\delta-e_{m+1})/(e_m-e_{m+1})\). The exact minimum worst expected paid cost is
\[
 \frac{\min_{i:r_i>0}(c_i/r_i)}{\mu}
 \left[t\,m(1-2e_m)+(1-t)(m+1)(1-2e_{m+1})\right].
\]
\end{theorem}
\begin{proof}
Under the equal prior on the matched pair the posterior odds are \(\lambda^s\). Conditional ordinary data have the same law given the receipts and can be simulated on the pair. The conditional minimum stop error is \(e_{|s|}\). Backward induction gives the first display, and reflecting bits makes the two type errors equal under symmetric tie-breaking. At a fixed reveal boundary \(j\), the probability of the wrong exit is \(e_j\) by the likelihood ratio. The expected number of reveals is \(j(1-2e_j)/\mu\) by summing the biased walk's conditional drifts to exit. For any finite-expected rule, summing those drifts in the two types gives mean reveal count \(\mathbb E[|s|(1-2e_{|s|})]/\mu\) under the equal prior. This follows by conditioning the signed drift sum on the stopped record and its posterior. Truncation justifies summation since bounded increments are dominated by the integrable stop count. The Bayes error is at least \(\mathbb E e_{|s|}\). The function \((1-2e)\log((1-e)/e)/\log\lambda\) is decreasing and convex on \(0<e\le1/2\); its second derivative is \(1/[e^2(1-e)^2\log\lambda]\). The secant through adjacent grid values \(e_m,e_{m+1}\) is a lower bound at every grid point. Thus no rule can beat the displayed interpolation if its error is at most \(\delta\). Each reveal costs at least \(\min_i c_i/r_i\) in expectation by summing conditional charges. A private lottery over the two boundaries with probabilities \(t,1-t\), using a minimum-price-per-reveal probe, attains equality. A zero boundary guesses fairly. At error zero, the overlapping finite paths exclude almost surely finite perfect classification. For equal prices the most revealing affordable probe repeated up to the move limit is optimal by thinning its reveals to simulate any other probe.
\end{proof}

\begin{proposition}[T30.A.4. Receipt controls]\label{flag:a4}
Take four table entries with one or three successes and four independent surcharge tickets, of which the responder selects zero or one. Use probes \((c_i,r_i)=(1,1/2),(3,1)\). At threshold \(1/2\) the positive linked risk is at most \(7/16\). Under indefinitely repeatable access the cost at \(\delta=1/10\) is \(32/5\); two moves cannot achieve this error at any budget.
\end{proposition}
\begin{proof}
Here \(q=3/4,\lambda=3,e_2=1/10\), and the cost per reveal is two. Substitute \(m=2,\mu=1/2\) into Theorem~\ref{flag:a3}. Even two complete reveals have optimal error \(1/4\) with a fair tie. Remove the protected receipts and matched ordinary laws give error at least \(1/2\). Alternatively allow a two-ticket proposal at the same positive type. The reference risks \(1/4\) and \(5/8\) of two matched-audit continuations straddle the threshold. Equal probe mechanics alone cannot resolve that change in the continuation class.
\end{proof}

\subsection{Finishing in the installation that produced the evidence}
\begin{theorem}[T30.U.1. Exact finite joint control]\label{flag:u1}
Fix a finite list of complete worlds, finite horizon \(H\), perfect recall, rational initial and transition rows, fixed intrinsic labels and an independent private controller seed. Include actual responder memory, phase cues, path resources and genuine event and service flags in the source state. Commands and stops must be executable on their complete reachable supports. The attainable vectors of bounded additive performance measures across worlds form a polytope. A rational linear program optimizes worst expected acquisition under affine constraints in every world. In particular it can constrain errors, coverage, actual violation and utility on issued runs, and the separate hard price. Uniform conditional guarantees require positive coverage on the worlds to which they apply.
\end{theorem}
\begin{proof}
Use the realization factors of Theorem~\ref{flag:p1} on all observable records with the conditional rows of the one world throughout an attempt. A terminal violation flag also covers risks earlier in the attempt. Expected scores are linear leaf sums, or sums of stage contributions. For selective risk bound \(\Pr(V\cap I)-\varepsilon\Pr(I)\) above by zero; for conditional utility bound \(\mathbb E[G1_I]-u\Pr(I)\) below by zero. Require positive exposure separately. A forbidden support outcome or hard overrun is a nonnegative marker with expectation zero in every world. Test each distinct attainable path price for an exact hard optimum. The projection of the compact sequence polytope is compact, giving attainment; linear infeasibility gives a rational separating certificate. Intrinsic label homogeneity, feasible audit-only identification and joint useful realization are three different tests. A missing physical correspondence or a computation that times out supplies none of their impossibility certificates.
\end{proof}

\begin{theorem}[T30.U.2. A simultaneous information state]\label{flag:u2}
Assign each listed world an artificial strictly positive rational mass \(p_w\). Keep the time, observed public register and posterior on the pair (world, hidden state), using the likelihood of the received record conditional on its imposed actions. When these quantities determine all relevant future laws, permissions and metrics, policies on this state realize every full-record policy's performance vector in every world. If each world's hidden start is fixed and the next hidden state is determined by its current state, command and received letter on support, the number of such states is at most
\[
 (H+1)|S|\prod_w\left(1+M_w\binom{H+r_w}{r_w}\right),
\]
where \(S\) is the public register, \(M_w\) bounds hidden states, and \(r_w\) counts distinct positive conditional letter likelihoods. Any initial reply must be counted in the horizon.
\end{theorem}
\begin{proof}
At equal posterior, time and public register, the conditional joint future given an action has the same kernel. Average the full-history controller's action choices with their actual occupation weights in the artificial mixture. Induction preserves the joint state, action and hidden-world occupation. Dividing that mixture's world contributions by \(p_w\) recovers every world coordinate, including zero-probability-of-breach checks. Positive prior mass everywhere is essential here. In the deterministic hidden-update case a surviving world needs just its hidden endpoint and counts for each distinct factor to compute the likelihood product. A dead-world symbol covers zero likelihood. The stars-and-bars bound counts vectors with total at most \(H\). No installation frequencies enter the conclusion; the expanded program and its exact arithmetic must still be paid for.
\end{proof}

\begin{theorem}[T30.U.3. Calibration and the issued denominator]\label{flag:u3}
Suppose external calibration contains the true completion with probability at least \(1-\eta\), and deployment conditional on the selected policy and calibration record obeys the checked kernel. On coverage let label error and whole-attempt risk be at most \(\delta,\varepsilon\). Their marginal bounds are \(\delta+(1-\delta)\eta\) and \(\varepsilon+(1-\varepsilon)\eta\); their union is bounded by \(\eta+(1-\eta)\min(1,\delta+\varepsilon)\). If instead the covered records satisfy a selected-risk cap \(\varepsilon\) and \(a=\Pr(I)>0\), then
\[
 \Pr(V\mid I)\le\varepsilon+(1-\varepsilon)\min(1,\eta/a).
\]
\end{theorem}
\begin{proof}
Condition on the calibration record. Give the failed-coverage event its worst possible value. In the selective case its issued mass is at most \(\min(a,\eta)\), so the numerator is bounded by \(\varepsilon a+(1-\varepsilon)\min(a,\eta)\). No independence of issuance and coverage is assumed.
\end{proof}

\begin{proposition}[T30.U.4. A routed stochastic service]\label{flag:u4}
Take a fixed sign, route and responder bit. A finite bit prelude of arbitrary length \(N\) depends on the responder and its preceding bit, with no sign information. Safe intake supplies a route-specific work receipt. One optional price-one dry read observes a bit of mean \(1/4\) or \(3/4\) by sign; the actual finalizer uses the same threshold on a fresh draw. On a positive verdict the literal route creates the required artifact. The responder may change its artifact note. The ordinary reference always finishes, with target cutoff \(2/5\). For error \(\delta\in[1/4,1/2)\), whole-attempt risk at most \(2/5\) and positive utility at least \(1-\delta\), the optimal worst expected paid price is \(2-4\delta\), with hard price one.
\end{proposition}
\begin{proof}
Matched opposite worlds agree on every pre-probe datum. Buying with probability \(h\) yields total variation at most \(h/2\). Error \(\delta\) forces \(h/2\ge1-2\delta\). Buy with this minimum probability, decide from the bit if bought, guess fairly otherwise and actually finish on positive verdicts. Worldwise errors equal \(\delta\); whole risk is \( (1-\delta)/4\) on positives and \(3\delta/4\) on negatives. Error below \(1/4\) is impossible. Any finite record of issuance with positive probability on the matched positive world also has positive probability on its negative counterpart, whose conditional risk is \(3/4\). A uniform selective cap \(2/5\) with positive service is therefore impossible here.

A stronger finite construction lets the note change the \emph{physical} event. In prelude phase take \(v=0\) initially, draw the next public bit with chance of one \((1+(m\mathbin{\mathrm{xor}} v))/3\) for fixed hidden bit \(m\), and retain the last bit \(v\). On a positive verdict the carried responder chooses a Boolean proposal using \(m,v\) and the verdict. The finalizer with base risk \(r\in\{1/4,3/4\}\) now has conditional risk \(r+(1-r)a/4\) for proposal \(a\). Check these rows at the actual reached state. They give positive cuts at most \(7/16\) and negative cuts at most \(13/16\). With threshold \(1/2\), price \(2/3\) buys error \(1/3\), positive safe usefulness at least \(3/8\), and whole risk at most \(7/24\). The pre-verdict lower bound is unchanged. Uniform low selective risk on both signs still fails. Theorem~\ref{adv:ug3} uses a different guarded family to obtain it.
\end{proof}

\subsection{Two limits at the decision boundary}
\begin{lemma}[Continuity at a threshold]\label{flag:continuity}
In the P3 interface let a positive threshold world \(w_0\) have the same exact tuple and complete startup law as negative worlds \(w_n\), with their entire diagnostic mean vectors converging to its vector. If positive issuance has probability at most \(\delta\) in every \(w_n\), then its probability at finite time in \(w_0\) is also at most \(\delta\).
\end{lemma}
\begin{proof}
At every fixed time limit \(H\), coupling the independent controller seed and each chosen bit on equal records bounds the discrepancy by \(H\max_j|m_j(w_n)-m_j(w_0)|\), which vanishes. Exact fields supply no further difference. Positive issuance by time \(H\) therefore has limiting probability at most \(\delta\) in \(w_0\). Increase \(H\) and use continuity of probability for increasing events. This rules out uniform error below \(1/2\) even for a procedure with a world-dependent almost surely finite stopping time at the boundary.
\end{proof}
\begin{lemma}[Attainable activation tradeoff in the finite one-read service]\label{flag:roc}
In a matched opposite pair of Theorem~\ref{flag:u4}, let \(q_+,q_-\) be their respective positive-issuance probabilities. Then \(q_-\ge\max(q_+/3,3q_+-2)\), with equality attainable for every \(q_+\) over the record tests alone.
\end{lemma}
\begin{proof}
The negative-to-positive likelihood ratio is one on skip and \(1/3\) or three for the two dry outcomes. Integrate the lower likelihood ratio on issuance and the upper ratio on its complement. To attain the boundary buy once, admit a randomized portion of zero outcomes until \(q_+=3/4\), then all zeros and a randomized portion of ones. A separate physical constraint can remove points on this curve. In the negative world joint violation is \(3q_-/4\), so requiring it to be at most \((2/5)q_-+\xi\) forces \(\xi\ge(7/20)\max(q_+/3,3q_+-2)\). This quantifies the exposure cost of even an approximate selective guarantee.
\end{proof}
